\section{OpenAI、常時稼働エージェント「dots」を発表}
OpenAIはDevDay 2026で、GPT-6 Astraを動力とする常時稼働エージェント「dots」を発表した。接続したプラグインを通じて多数のアプリを横断し、独自のクラウドコンピュータ上で24時間バックグラウンド作業を行うと説明している。Pro、Business Premium、Enterpriseの対象市場でChatGPT web/mobile/desktopへ提供するとしている。
\dailyxsource{OpenAI (@OpenAI)}{https://x.com/OpenAI/status/2104984504133918973}
\dailyxnote{対象地域、料金、権限モデル、ログ保持の詳細はGrok収集では未精読。}

\section{GPT-6.1 Solを公開}
OpenAIはGPT-6.1 Solを「near-Astra intelligence for a fifth of the price」と位置づけて公開した。Plus以上のChatGPT WorkとCodexで当日から利用可能とし、cached inputは100万token当たり0.10ドル、標準入力比95\%安などを公式スレッドで主張している。Astraとの比較条件や公式ベンチマーク表の詳細はGrok収集では未突合である。
\dailyxsource{OpenAI (@OpenAI)}{https://x.com/OpenAI/status/2104986129686741046}

\section{Astra UltrafastとPro 500を発表}
OpenAIはGPT-6 Astra向けのプレミアム速度ティア「Ultrafast」を発表し、Codexで最大8倍・300 tokens/sec、APIで最大6倍と主張した。Codex、ChatGPT Work、APIで当日提供し、GPT-6.1 Solへの対応は近日としている。併せてPro 500を導入し、Plus比25倍の上限と説明したが、月額などの詳細は公式投稿本文では確認されていない。
\dailyxsource{OpenAI (@OpenAI)}{https://x.com/OpenAI/status/2104993966043320759}

\section{Codex Security Cloudを強化}
OpenAIはCodex Security Cloudの大型更新を発表し、サイバー能力モデル「Daybreak Blue」へのアクセスを標準で含めるとした。GitHubリポジトリ全体のスキャン、新規コミットの継続レビュー、発見事項の調査・重複排除、修正案準備を、端末を閉じている間も行うと説明している。専用ドキュメントURLはGrok収集では確認されていない。
\dailyxsource{OpenAI (@OpenAI)}{https://x.com/OpenAI/status/2104987422308335828}

\section{Artificial AnalysisがGPT-6.1 Solを評価}
Artificial AnalysisはGPT-6.1 Solについて、Intelligence IndexでGPT-6 Astraに1ポイント差、タスク当たりコストはAstraの4分の1未満と報告した。Terminal-Bench 4.0やHLEなど複数指標の改善も示しているが、これらはAA独自の評価結果である。OpenAI公式価格との完全な突合やeffort設定の対応は未確認である。
\dailyxsource{Artificial Analysis (@ArtificialAnlys)}{https://x.com/ArtificialAnlys/status/2105025585332605357}
\dailyxnote{Grok intakeではLikely。数値は独立評価機関側の測定。}

\section{ホワイトハウスでSuper Intelligence会合}
White House Press Poolは、ホワイトハウスでSuper Intelligence会合が開かれ、トランプ大統領が署名直後の文書を「almost like a constitution」「morally binding」と表現したと速報した。自己規制や約10人規模の監視委員会構想への言及も報じられたが、協定本文や署名者の公式一覧はGrok収集では確認されていない。
\dailyxsource{White House Press Pool Reports (@WHPressPool)}{https://x.com/WHPressPool/status/2105030810910744717}
\dailyxnote{Grok intakeではLikely。4層監督や署名企業の整理は二次投稿由来。}

\section{America.govを政府サービスの単一入口として発表}
ホワイトハウスはAmerica.govを連邦政府サービスの新しい全国フロントドアとして公式発表し、ファクトシートURLも提示した。生成AIの実装分担については、Grok回答がGeminiとGrokを技術パートナーと説明しているが、政府ファクトシート本文との一致はGrok収集では未検証である。サイト発表とAIベンダー情報を分けて扱う必要がある。
\dailyxsource{The White House (@WhiteHouse)}{https://x.com/WhiteHouse/status/2104967240059687393}
\dailyxnote{サイト発表はConfirmed、各社の実装分担詳細はUnverified。}

\section{Anthropic Interviewer調査を実施}
AnthropicはAnthropic Interviewerによる新たな質的調査を9月29日から10月6日まで実施すると発表した。前回は81,000人が回答したとし、今回は回答を公開するかを任意で選べるとしている。ClaudeおよびClaude CodeのFree、Pro、Maxユーザーが対象で、公開回答を研究や意思決定に活用すると説明している。
\dailyxsource{Anthropic (@AnthropicAI)}{https://x.com/AnthropicAI/status/2104982629884063840}

\section{GLM-5.3の高度サイバー能力が議論に}
Anthropicの研究ページ「GLM-5.3 and the spread of advanced cyber capabilities」が観測窓終盤に広く共有された。二次要約では、オープンウェイトのGLM-5.3が閉じたMythos Previewに近いExploitBench成績を示したとされる。本文数値はGrok収集で未精読であり、オープンモデルの脅威評価を巡って安全側とオープンソース擁護側の反応が分かれた。
\dailyxsource{Ethan Mollick (@emollick)}{https://x.com/emollick/status/2105052282354303358}
\dailyxnote{Grok intakeではLikely。研究ページの詳細数値は二次要約段階。}

\section{frontier AI訓練のsafety caseを公開}
OpenAIは「Towards safety cases for frontier AI training」を公開し、frontier RL訓練ランの安全性を、技術的対策、運用管理、インシデント調査などを組み合わせたsafety caseとして文書化する枠組みを示した。Greg Brockmanも現行の学びを反映した実務ガイドラインとして紹介したが、必須の内部プロセスなのか公開ガイドラインなのかは本文確認が必要である。
\dailyxsource{OpenAI (@OpenAI)}{https://x.com/OpenAI/status/2104815409522483470}
